% Ad Familiares 15.20 (ad-familiares-15-20)
% Source: https://raw.githubusercontent.com/PerseusDL/canonical-latinLit/master/data/phi0474/phi056/phi0474.phi056.perseus-lat1.xml
% Fetched by scripts/fetch_latin.py

% Dateline (Perseus): Scr. Romae (?) m. Apr. 710 (44) .

\ciceroLetterOpener{M. CICERO S. D. C. TREBONIO}

\ciceroSection{1} ' oratorem ' meum (sic enim inscripsi) Sabino tuo commendavi. natio me hominis impulit ut ei recte putarem ; nisi forte candidatorum licentia hic quoque usus hoc subito cognomen arripuit; etsi modestus eius vultus sermoque constans habere quiddam a Curibus videbatur. sed de Sabino satis.

\ciceroSection{2} tu , mi Treboni, quoniam ad amorem meum aliquantum olei discedens addidisti, quo tolerabilius feramus igniculum desideri tui, crebris nos litteris appellato, atque ita, si idem fiet a nobis. quamquam duae causae sunt cur tu frequentior in isto officio esse debeas quam nos, primum quod olim solebant qui Romae erant ad provincialis amicos de re p. scribere, nunc tu nobis scribas oportet (res enim publica istic est) ; deinde quod nos aliis officiis tibi absenti satis facere possumus, tu nobis nisi litteris non video qua re alia satis facere possis.

\ciceroSection{3} sed cetera scribes ad nos postea ; nunc haec primo cupio cognoscere iter tuum cuius modi sit, ubi Brutum nostrum videris, quam diu simul fueris ; deinde cum processeris longius, de bellicis rebus, de toto negotio, ut existimare possumus quo statu simus. ego tantum me scire putabo quantum ex tuis litteris habebo cognitum. cura ut valeas meque ames amore illo tuo singulari.
